% A Finite Obstruction to Heisenberg Carrier Selection
% from Admissibility Constraints

\documentclass[11pt,a4paper]{article}

\usepackage[T1]{fontenc}
\usepackage[utf8]{inputenc}
\usepackage{lmodern}
\input{glyphtounicode}
\pdfgentounicode=1
\usepackage{amsmath,amssymb,amsthm}
\usepackage{enumerate}
\usepackage{booktabs}
\usepackage{array}
\usepackage{microtype}
\usepackage[hidelinks]{hyperref}
\usepackage{doi}

\newtheorem{theorem}{Theorem}[section]
\newtheorem{lemma}[theorem]{Lemma}
\newtheorem{proposition}[theorem]{Proposition}
\newtheorem{corollary}[theorem]{Corollary}
\newtheorem{definition}[theorem]{Definition}
\newtheorem{remark}[theorem]{Remark}

\newcommand{\Heis}{\mathrm{Heis}_3}
\newcommand{\Sym}{\mathfrak{S}}

\title{A Finite Obstruction to Heisenberg Carrier Selection\\
from Admissibility Constraints}

\author{J\'er\^ome Beau\\
\small Independent Researcher\\
\small \texttt{jerome.beau@cosmochrony.org}}

\date{2026}

\begin{document}

  \maketitle

  \begin{abstract}
    We test whether the admissibility axioms of the Cosmochrony framework,
    together with Born--Infeld parity, select a finite Heisenberg group and its
    irreducible carrier.
    The load-bearing algebraic implication is that a finite faithful irreducible
    unitary carrier, equipped with a minimally generating pair exchanged by an
    involution and having a non-trivial commutator, must have a central commutator.
    We disprove that implication by an explicit six-element countermodel.
    Let $G_n=\Sym_3$, $X=(12)$, $Y=(23)$, and let $\sigma$ be conjugation by $(13)$.
    The involution exchanges $X$ and $Y$, the pair minimally generates $G_n$, and
    the standard action on the two-dimensional zero-sum subspace of
    $\mathbb{C}^3$ is faithful, unitary, and irreducible.
    Nevertheless, $[X,Y]$ is a non-trivial three-cycle and is not central.
    We also show that irreducibility does not force the centre to have prime
    order: a central subgroup acts through one character, not through several
    non-zero character eigenspaces.
    Stone--von Neumann remains a valid conditional endpoint once a finite
    Heisenberg group and a non-trivial central character are supplied.
    The associated Weil action is a further representation of the symplectic
    automorphism group, not the Heisenberg representation itself.
    \textbf{Interpretive status.}
    The result obstructs the unique-selection bridge but leaves intact the
    mathematics proved internally on a supplied Heisenberg carrier or associated
    Weil module.
  \end{abstract}

  \paragraph{Keywords.}
    finite countermodel; carrier selection; Heisenberg group; admissibility;
    irreducible representation; Schur's lemma; nilpotence; Stone--von Neumann theorem.

    \clearpage
    \tableofcontents
    \clearpage

  \section{Introduction}
  \label{sec:intro}

  The Cosmochrony programme~\cite{Cosmochrony} proposes a relational foundation
  in which admissible transitions precede effective quantum and geometric
  descriptions.
  The companion Foundation paper~\cite{Foundation} formulates four axioms and
  identifies the admissible fibre $F_n$ with the carrier used in a finite
  Heisenberg--Weil construction.
  The purpose of the present paper is to determine whether the algebraic
  selection step follows from the stated inputs.

  The main result, Theorem~\ref{thm:main}, is:

  \begin{quote}
    \textit{The finite, faithful, irreducible, unitary, minimal-pair contract
    extracted before the centrality step does not select a Heisenberg carrier.}
  \end{quote}

  \noindent
  The obstruction is exact and finite.
  The group $\Sym_3$ satisfies the full contract while its non-trivial
  commutator is not central.
  A separate application of Schur's lemma blocks the prime-order step.
  The result distinguishes what follows from the admissibility inputs from what
  follows only after a finite Heisenberg group has been supplied.

  \subsection*{Relation to companion papers}

    This paper is a structural companion to~\cite{Foundation}, which states
    Axioms A1--A4 and derives their consequences.
    The non-abelian and quaternionic constructions discussed in
    Remark~\ref{rem:nonabelian} are developed in~\cite{O23}.
    The Born--Infeld parity involution $\sigma\colon c \mapsto q - c$ is derived
    in~\cite{BornInfeld}.
    The non-injectivity of $\Pi_n$ as a structural necessity is proved
    in~\cite{ENI}.
    None of these companion results supplies the missing class-two or central
    extension hypothesis used in the selection step.

  \subsection{Relation to neighbouring constructions}
  \label{sec:neighbouring}

    The obstruction does not say that a finite Weyl--Heisenberg structure cannot
    be selected from explicit structural inputs.
    Mackey's imprimitivity programme gives a precise sufficient route.
    For a finite Abelian configuration group $A$, an irreducible transitive
    system of imprimitivity consists of a unitary representation $U$ and a
    projection-valued localisation $E$ satisfying
    \[
      U(a)E(S)U(a)^{-1}=E(a^{-1}S).
    \]
    These data determine the finite quantum kinematics up to unitary equivalence
    and yield the associated discrete Weyl system~\cite{StovicekTolar1984,Tolar2014}.

    A complementary mathematical route starts with a finite symplectic vector
    space $(V,\omega)$ over a field of odd characteristic.
    Its non-degenerate alternating form defines a Heisenberg central extension;
    Stone--von Neumann then fixes the irreducible Heisenberg representation, and
    symplectic functoriality produces the associated Weil
    representation~\cite{GurevichHadani2009}.

    Both routes supply more typed structure than the audit contract used here.
    The contract specifies neither a finite Abelian configuration space with
    covariant localisation nor a symplectic vector space with a non-degenerate
    form.
    These neighbouring constructions therefore show what a positive selection
    theorem can use without repairing the failed implication from minimal
    generation alone.

  \subsection*{Organisation}

    Section~\ref{sec:setup} recalls the axiomatic setup and fixes notation.
    Section~\ref{sec:lemma} states the representation-theoretic reading of
    admissible non-factorisability.
    Section~\ref{sec:main} gives the countermodel, the Schur correction, and the
    conditional Stone--von Neumann endpoint.
    Section~\ref{sec:corollaries} records the consequences available only after
    a carrier is supplied.
    Section~\ref{sec:discussion} states the resulting epistemic boundary.

    \clearpage
  \section{Setup and Notation}
  \label{sec:setup}

  We recall the axioms and fix the notation used throughout.
  The full axiomatic development is given in~\cite{Foundation}.

  \subsection{The four axioms}

    \begin{description}
      \item[\textbf{A1.}] \textbf{Local projective admissibility.}
      For any observable state $O_{n-1}$, there exists a non-empty set $F_n$ of
      admissible successor directions.
      The set $F_n$ is defined by internal structural constraints (Born--Infeld
      bound, spectral envelope, symmetry requirements) and does not presuppose
      an ambient configuration space.

      \item[\textbf{A2.}] \textbf{Structural non-injectivity.}
      Distinct admissible directions from $O_{n-1}$ may lead to the same resolved
      successor state $O_n$.

      \item[\textbf{A3.}] \textbf{Admissibility as non-premature selection.}
      An admissible transition preserves the full multiplicity of open directions
      until projection locking occurs.
      It does not select among them prematurely, that is, before the
      Born--Infeld saturation threshold is reached.

      \item[\textbf{A4.}] \textbf{Projection locking.}
      Resolution occurs only when continuous Born--Infeld saturation meets the
      discrete support of the observable image of the projection.
      In the spectral admissibility realisation, this support is the shell
      decomposition of the observable space.
    \end{description}

  \subsection{The admissible fibre and its generators}

    For the finite carrier-selection problem, let $F_n$ be a finite-dimensional
    complex vector space carrying a faithful unitary action of a finite group
    $G_n$.
    These are explicit inputs to the algebraic contract tested below.
    An amplitude bound alone does not establish finite dimensionality, so the
    latter is not counted as a consequence of the Born--Infeld bound.

    By A1, the admissible fibre $F_n$ is a non-empty set of admissible directions
    from $O_{n-1}$.
    It carries a group of symmetry transformations, whose generators are constrained
    by the Born--Infeld bound and the spectral envelope.
    We denote the symmetry group of $F_n$ by $G_n$; its action on $F_n$ is by
    definition faithful.

    The selection argument uses an involutive automorphism $\sigma$ of $G_n$ and
    a generating pair $(X,\sigma(X))$ exchanged by that involution.
    It further assumes that the pair is minimal, meaning that neither singleton
    generates $G_n$.
    The finite character map $c\mapsto q-c$ is available only after a cyclic
    centre of order $q$ and its character labelling have been supplied.
    Continuum Born--Infeld parity~\cite{BornInfeld} does not by itself construct
    those finite data.

    \begin{definition}[Stable decomposition]
      \label{def:stable}
      A direct sum decomposition $F_n = F^{(1)} \oplus F^{(2)}$ is called
      \emph{stable} (under admissible generators) if both $F^{(1)}$ and $F^{(2)}$
      are invariant under the action of every generator of $G_n$.
    \end{definition}

    \begin{definition}[Admissible resolution]
      \label{def:resolution}
      A direction $f \in F_n$ is \emph{admissibly resolved} if there exists a
      stable decomposition $F_n = F^{(1)} \oplus F^{(2)}$ such that
      $f \in F^{(1)}$ or $f \in F^{(2)}$, i.e.\ $f$ is assigned to a definite
      sector prior to any Born--Infeld saturation event.
    \end{definition}

  \subsection{The Heisenberg representation and associated Weil action}

    For a prime $q$ and a non-zero character $c \in (\mathbb{Z}/q\mathbb{Z})^\times$,
    the discrete Heisenberg group
    $\mathrm{Heis}_3(\mathbb{Z}/q\mathbb{Z}) = \{(a,b,\gamma) : a,b,\gamma \in \mathbb{Z}/q\mathbb{Z}\}$
    with multiplication
    $(a,b,\gamma)(a',b',\gamma') = (a+a',\, b+b',\, \gamma+\gamma'+ab')$
    admits an irreducible representation of dimension $q$, the Heisenberg or
    Schr\"odinger representation $\pi_c$, acting on the carrier
    $\mathcal{H}_c=L^2(\mathbb{Z}/q\mathbb{Z})\cong\mathbb{C}^q$ by
    \begin{equation}
      \bigl(\pi_c(a,b,\gamma)\,f\bigr)(x)
      \;=\; e^{2\pi i c(\gamma + bx)/q}\,f(x - a).
      \label{eq:heisenberg-representation}
    \end{equation}
    The central element $(0,0,1)$ acts as multiplication by $e^{2\pi ic/q} \neq 1$
    (because $c \neq 0$), so the central character is non-trivial.
    The Weil representation is the associated action of the symplectic
    automorphism group on the same carrier, constructed through intertwiners of
    the Heisenberg representation~\cite{GurevichHadani2009}.
    Stone--von Neumann identifies $\pi_c$; it does not by itself supply that
    symplectic action.

    \clearpage
  \section{The Admissible Non-Factorisability Lemma}
  \label{sec:lemma}

  The following lemma records the representation-theoretic reading used for
  Axiom A3.
  Its status is conditional on identifying an invariant subspace with a
  prematurely resolved sector.

  \begin{lemma}[Admissible non-factorisability]
    \label{lem:non-factorisation}
    Under the representation-theoretic reading of A3, every stable decomposition
    of $F_n$ is trivial, i.e.\ one of $F^{(1)}$ and $F^{(2)}$ must be $\{0\}$.
    In particular, the action of $G_n$ on $F_n$ is irreducible.
  \end{lemma}

  \begin{proof}
    Suppose $F_n = F^{(1)} \oplus F^{(2)}$ is a non-trivial stable decomposition,
    with $F^{(1)} \neq \{0\}$ and $F^{(2)} \neq \{0\}$.
    By Definition~\ref{def:stable}, both summands are invariant under every
    generator of $G_n$.
    Since the generators of $G_n$ are the constitutive elements of the
    admissibility structure of $F_n$ (A1), the decomposition is distinguished
    by the admissibility structure itself: the sector to which a direction $f$
    belongs ($F^{(1)}$ or $F^{(2)}$) is determined by an invariant internal to $F_n$.

    This constitutes an admissible resolution in the sense of
    Definition~\ref{def:resolution}: each open direction carries a definite
    sectorial label, accessible within the admissibility structure of $F_n$,
    prior to any Born--Infeld saturation event.

    Under the stated representation-theoretic reading, A3 forbids premature
    selection among open directions: the transition must remain structurally
    neutral among all elements of $F_n$ until projection locking forces
    resolution.
    An admissible resolution prior to locking is what that reading prohibits.
    Therefore the non-trivial stable decomposition is forbidden.

    Irreducibility follows immediately: any proper $G_n$-invariant subspace
    $W \subsetneq F_n$ with $W \neq \{0\}$ would yield a non-trivial stable
    decomposition $F_n = W \oplus W^\perp$ (because $G_n$ acts unitarily, the
    orthogonal complement of a $G_n$-invariant subspace is also $G_n$-invariant).
    This contradicts the non-factorisability just proved.
  \end{proof}

  \begin{remark}[Status of Lemma~\ref{lem:non-factorisation}]
    \label{rem:restatement}
    The content of Lemma~\ref{lem:non-factorisation} may be expressed as:
    \begin{equation}
      \begin{aligned}
        &\text{proto-state genuinely non-resolved}\\
        &\qquad\Longleftrightarrow
        \text{ no admissible invariant-sector decomposition of }F_n.
      \end{aligned}
    \end{equation}
    The equivalence is a dictionary between the physical language of resolution
    and the algebraic language of invariant sectors.
    The lemma is therefore conditional on that dictionary; it is not an
    independent derivation of irreducibility from the other axioms.
  \end{remark}

  \begin{remark}[Unitarity of the $G_n$ action]
    \label{rem:unitarity}
    The action of $G_n$ on $F_n$ is assumed unitary in the finite algebraic
    contract.
    A bound expressed through mode amplitudes does not alone prove that every
    admissible symmetry preserves the corresponding Hermitian norm.
    Unitarity is used in the proof above to ensure that the orthogonal complement
    of an invariant subspace is again invariant.
  \end{remark}

  \section{The Carrier-Selection Obstruction}
  \label{sec:main}

  \subsection{Audit contract}

    Before any centrality or classification input, the finite selection argument
    has the following algebraic data:
    \begin{enumerate}
      \item a finite-dimensional complex carrier $F_n$;
      \item a faithful irreducible unitary action of a finite group $G_n$;
      \item a generating pair $X,\sigma(X)$, with no proper subpair generating $G_n$;
      \item an involutive automorphism $\sigma$ exchanging the two generators;
      \item a non-trivial group commutator $[X,\sigma(X)]$.
    \end{enumerate}
    The question is whether these data force the commutator to be central.
    Permutations below are composed from right to left, and
    $[a,b]=aba^{-1}b^{-1}$.

  \subsection{The six-element countermodel}

    \begin{theorem}[Finite carrier-selection obstruction]
      \label{thm:main}
      Let
      \[
        G_n=\Sym_3,\qquad X=(12),\qquad Y=(23),\qquad \tau=(13),
      \]
      and define $\sigma(g)=\tau g\tau^{-1}$.
      Let $G_n$ act by coordinate permutation on
      \[
        F_n=\bigl\{(z_1,z_2,z_3)\in\mathbb{C}^3:
        z_1+z_2+z_3=0\bigr\}.
      \]
      Then:
      \begin{enumerate}
        \item $\sigma$ is an involutive automorphism and exchanges $X$ and $Y$;
        \item $X$ and $Y$ form a minimal generating pair for $G_n$;
        \item the representation on $F_n$ is two-dimensional, faithful, unitary,
        and irreducible;
        \item $[X,Y]\neq 1$;
        \item $[X,Y]\notin Z(G_n)$.
      \end{enumerate}
      Consequently, the audit contract does not force class-two nilpotence, a
      finite Heisenberg group, or its irreducible carrier.
    \end{theorem}

    \begin{proof}
      Because $\tau^2=1$, conjugation by $\tau$ satisfies
      $\sigma^2=\mathrm{id}$.
      Conjugation relabels the entries of a transposition, so
      \[
        \sigma((12))=(\tau(1)\,\tau(2))=(32)=(23)=Y,
      \]
      and similarly $\sigma(Y)=X$.

      The elements $X$ and $Y$ generate $\Sym_3$ because $XY=(123)$.
      Neither singleton generates $G_n$, because each generates a subgroup of
      order two.

      Coordinate permutations preserve the standard Hermitian product on
      $\mathbb{C}^3$, and the zero-sum subspace $F_n$ is invariant.
      The restricted representation is therefore unitary and two-dimensional.
      If $g\in\Sym_3$ acts trivially on $F_n$, it fixes both $e_1-e_2$ and
      $e_2-e_3$, hence fixes every coordinate position and is the identity.
      The representation is faithful.

      Suppose that $F_n$ contains a $G_n$-invariant complex line $L$.
      The involutions $X$ and $Y$ act on $L$ through scalars in $\{1,-1\}$.
      Because they are conjugate, the scalars are equal.
      If both act as $+1$, a vector of $L$ has three equal coordinates and the
      zero-sum condition makes it zero.
      If both act as $-1$, then
      \[
        L\subseteq
        \mathbb{C}(e_1-e_2)\cap\mathbb{C}(e_2-e_3)=\{0\}.
      \]
      Both alternatives contradict the existence of $L$.
      The two-dimensional representation is irreducible.

      Finally, $X^{-1}=X$ and $Y^{-1}=Y$, so
      \[
        [X,Y]=XYXY=(XY)^2=(132)\neq 1.
      \]
      It is not central because
      \[
        X(132)X^{-1}=(123)\neq(132).
      \]
    \end{proof}

  \subsection{What non-factorisability does establish}

    \begin{proposition}
      \label{prop:commutator}
      Let $\dim F_n>1$.
      If the unitary action of $G_n=\langle X,\sigma(X)\rangle$ on $F_n$ is
      irreducible, then $X$ and $\sigma(X)$ do not commute as operators.
    \end{proposition}

    \begin{proof}
      If the two unitary generators commute, the spectral theorem gives a
      simultaneous eigenspace decomposition.
      Every simultaneous eigenspace is invariant under both generators and
      therefore under $G_n$.
      Irreducibility makes the decomposition consist of a single eigenspace, so
      both generators act as scalars.
      The entire representation is then a direct sum of one-dimensional
      representations, contradicting irreducibility when $\dim F_n>1$.
    \end{proof}

    Proposition~\ref{prop:commutator} is an elementary consequence of the
    algebraic contract once the non-trivial dimension and irreducibility inputs
    are explicit.
    It does not constrain higher commutators.

    \begin{proposition}[Failure of the centrality implication]
      \label{prop:central}
      Minimal generation by an involutively exchanged pair, faithful
      irreducibility, unitarity, and a non-trivial commutator do not imply that
      the commutator is central.
    \end{proposition}

    \begin{proof}
      Theorem~\ref{thm:main} satisfies every premise and violates the conclusion.
    \end{proof}

    \begin{remark}[Nilpotency class two is an additional input]
      \label{rem:nilpotency}
      A relation such as $[[X,\sigma(X)],X]=1$ is not a new generator.
      It is a relation among elements already generated by the pair.
      Minimality of a generating set therefore does not bound the lower central
      series.
      Class-two nilpotence must be assumed or derived independently.
    \end{remark}

  \subsection{A sufficient symplectic input}

    \begin{proposition}[Heisenberg group from supplied symplectic data]
      \label{prop:heisenberg-algebra}
      Let $q$ be an odd prime.
      Let $(V,\omega)$ be a two-dimensional symplectic vector space over
      $\mathbb{F}_q$.
      On the set $H(V)=V\times\mathbb{F}_q$, define
      \begin{equation}
        (v,z)(v',z')
        =\bigl(v+v',z+z'+\tfrac12\omega(v,v')\bigr).
        \label{eq:heisenberg-presentation}
      \end{equation}
      Then $H(V)$ is a group of order $q^3$ with
      \[
        Z(H(V))=\{0\}\times\mathbb{F}_q,\qquad
        [(v,z),(v',z')]=(0,\omega(v,v')).
      \]
      It is isomorphic to $\Heis(\mathbb{F}_q)$.
    \end{proposition}

    \begin{proof}
      Bilinearity of $\omega$ gives associativity, and $(0,0)$ is the identity.
      The inverse of $(v,z)$ is $(-v,-z)$ because $\omega(v,v)=0$.
      Direct multiplication gives the stated commutator.
      Non-degeneracy implies that an element $(v,z)$ commutes with all of $H(V)$
      only when $v=0$, which gives the centre.
      Choosing a symplectic basis of $V$ identifies the multiplication law with
      the standard finite Heisenberg presentation~\cite{GurevichHadani2009}.
    \end{proof}

    Proposition~\ref{prop:heisenberg-algebra} is a construction from supplied
    symplectic data.
    The audit contract does not select $V$ or $\omega$.

  \subsection{The prime-order step}

    \begin{lemma}[Central subgroups on an irreducible carrier]
      \label{lem:central-character}
      Let $V$ be a finite-dimensional irreducible complex representation of a
      group $G$, and let $H\leq Z(G)$.
      Then $H$ acts on $V$ through a single character.
      Equivalently, in
      \[
        V=\bigoplus_{\chi\in\widehat H}V_\chi,
      \]
      at most one summand is non-zero.
    \end{lemma}

    \begin{proof}
      For $h\in H$, any eigenspace of the complex operator representing $h$ is
      $G$-invariant because $h$ is central.
      Irreducibility makes that eigenspace all of $V$.
      Thus each $h$ acts as a scalar, and the resulting scalar map is one
      character of $H$.
    \end{proof}

    \begin{proposition}[No prime-order conclusion from irreducibility]
      \label{prop:prime}
      The existence of a proper non-trivial subgroup of $Z(G_n)$ does not
      contradict irreducibility of a complex $G_n$-representation.
      Irreducibility alone therefore does not imply that the centre, or a
      specified central cyclic subgroup, has prime order.
    \end{proposition}

    \begin{proof}
      By Lemma~\ref{lem:central-character}, the whole carrier may transform
      through one non-trivial character of the subgroup.
      The number of available characters does not imply that several occur.
      The resulting character decomposition has one non-zero summand and is
      therefore not a non-trivial stable decomposition.
    \end{proof}

  \subsection{Conditional Stone--von Neumann endpoint}

    \begin{proposition}
      \label{prop:vrho}
      Suppose that $G_n\cong\Heis(\mathbb{Z}/q\mathbb{Z})$ and that $F_n$ is an
      irreducible complex representation with a specified non-trivial central
      character.
      Then $F_n$ is unitarily equivalent to the corresponding Heisenberg
      representation on $L^2(\mathbb{Z}/q\mathbb{Z})$.
    \end{proposition}

    \begin{proof}
      This is the finite Stone--von Neumann theorem~\cite{GurevichHadani2009}.
      For each non-trivial central character, the irreducible representation is
      unique up to unitary equivalence and is realised
      by~\eqref{eq:heisenberg-representation}.
    \end{proof}

    The finite character involution $c\mapsto q-c$ presupposes a cyclic centre of
    order $q$ and its character labelling.
    It cannot select the finite carrier before those objects exist.

    \begin{remark}[Typed derivation chain]
      \label{rem:chain}
      The available chain has four differently typed edges:
      \begin{itemize}
        \item the representation-theoretic reading of A3 implies irreducibility
        conditionally on the resolution dictionary;
        \item irreducibility, $\dim F_n>1$, and generation by the pair prove
        $[X,\sigma(X)]\neq 1$;
        \item a non-trivial commutator does not imply centrality, by
        Theorem~\ref{thm:main};
        \item a supplied Heisenberg group and central character imply
        $F_n\simeq\mathcal{H}_c$ as a Heisenberg representation, by
        Stone--von Neumann.
      \end{itemize}
      The missing edge is carrier selection, not the internal representation
      theory of a supplied carrier.
    \end{remark}

  \clearpage
  \section{Conditional Consequences}
  \label{sec:corollaries}

  \subsection{Uncertainty statements}

    \begin{corollary}[Conditional Robertson bound]
      \label{cor:uncertainty}
      Let $A$ and $B$ be self-adjoint observables on a supplied Hilbert carrier.
      For every normalised state $\psi$,
      \[
        \Delta_\psi A\,\Delta_\psi B
        \geq \frac{1}{2}\bigl|\langle\psi,[A,B]\psi\rangle\bigr|.
      \]
    \end{corollary}

    \begin{proof}
      This is the Robertson inequality, applied after a Hilbert carrier and
      self-adjoint observables have been specified.
    \end{proof}

    \begin{remark}[Interpretive boundary]
      \label{rem:uncertainty-interpretation}
      A non-zero operator commutator does not make its expectation strictly
      positive in every state, and the finite Heisenberg generators are unitary
      rather than automatically self-adjoint observables.
      Identifying their uncertainty relation with a physical measurement
      principle therefore requires an observable dictionary in addition to the
      carrier.
    \end{remark}

  \subsection{Finite Weyl quantisation}

    \begin{corollary}[Conditional finite Weyl system]
      \label{cor:quantisation}
      If the finite Heisenberg group, its non-degenerate symplectic pairing, and
      a non-trivial central character are supplied, then the carrier
      $\mathcal{H}_c$ supports the corresponding finite Weyl system.
    \end{corollary}

    \begin{proof}
      Proposition~\ref{prop:heisenberg-algebra} identifies the supplied central
      extension, and Proposition~\ref{prop:vrho} identifies its irreducible
      carrier.
      This is the mathematical quantisation of the supplied finite symplectic
      data.
    \end{proof}

  \subsection{Relation to the O-series}

    \begin{remark}[Scope of non-abelian support]
      \label{rem:nonabelian}
      The quaternionic construction in~\cite{O23} and the finite countermodel in
      Theorem~\ref{thm:main} answer different questions.
      A non-abelian neutral sector does not by itself determine the ambient
      symmetry group, its nilpotency class, or a finite Heisenberg presentation.
      Results proved on a chosen Heisenberg graph remain valid on that graph but
      are conditional when read as consequences of carrier selection.
    \end{remark}

  \clearpage
  \section{Discussion}
  \label{sec:discussion}

  \subsection{Established result}

    Theorem~\ref{thm:main} is a complete finite obstruction to the claimed
    selection implication.
    Its premises and failure of centrality are checked inside the six-element
    group $\Sym_3$.
    Proposition~\ref{prop:prime} independently shows that irreducibility supplies
    no prime-order conclusion.

    The result is local to the bridge
    \[
      \text{admissibility contract}
      \longrightarrow
      \text{finite Heisenberg carrier}.
    \]
    It does not identify $\Sym_3$ as a physical carrier.

  \subsection{What survives}

    The obstruction leaves intact:
    \begin{itemize}
      \item the elementary non-commutation result under the explicit hypotheses
      of Proposition~\ref{prop:commutator};
      \item the classification of a finite Heisenberg group after the required
      central-extension data are supplied;
      \item the finite Stone--von Neumann theorem for a specified non-trivial
      central character;
      \item results in companion programmes whose statements are internal to a
      chosen Heisenberg graph, Heisenberg carrier, or associated Weil module.
    \end{itemize}

  \subsection{Open selection problem}

    Excluding the $\Sym_3$ countermodel requires an additional input such as
    class-two nilpotence, a specified central extension, or an equivalent
    non-degenerate symplectic commutator law.
    Any such input must be motivated or derived independently.
    Calling higher commutators ``excess structure'' does not provide an algebraic
    implication.

    The large-$q$ limit, in which a supplied finite Heisenberg carrier and its
    associated Weil action may approach a real Schr\"odinger and metaplectic
    construction, remains a conditional direction of the spectral geometry
    programme~\cite{Cosmochrony}.
    It cannot resolve the finite carrier-selection edge retroactively.

  \subsection*{Interpretive outlook}

    \textbf{Interpretive status.}
    The paper separates two questions: what follows mathematically on a
    Heisenberg carrier, and why that carrier should describe the admissible
    fibre.
    The first question retains a substantial body of conditional mathematics.
    The second is an open selection problem.
    Keeping the two levels distinct permits continued exploration of the carrier
    without presenting the exploratory choice as an axiomatic necessity.

  \section*{Acknowledgements}

  The author acknowledges the use of large language models as a supportive tool
  for refining language, structure, and internal consistency during the development
  of this manuscript.
  All conceptual contributions, theoretical choices, and interpretations remain
  the sole responsibility of the author.

  \bibliographystyle{plain}
      \bibliography{cosmochrony-bibliography}

\end{document}
